\documentclass[12pt,letterpaper]{article}

% ─── Packages ────────────────────────────────────────────────────────────────
\usepackage[margin=1in]{geometry}
\usepackage{amsmath,amssymb,amsthm}
\usepackage{mathtools}
\usepackage{bm}
\usepackage{hyperref}
\usepackage{xcolor}
\usepackage{listings}
\usepackage{booktabs}
\usepackage{array}
\usepackage{longtable}
\usepackage{enumitem}
\usepackage{titlesec}
\usepackage{fancyhdr}
\usepackage{graphicx}
\usepackage{microtype}
\usepackage{abstract}
\usepackage{tcolorbox}
\tcbuselibrary{breakable,skins}
\usepackage{mdframed}
\usepackage{setspace}
\usepackage{natbib}

% ─── Colours ─────────────────────────────────────────────────────────────────
\definecolor{mtblue}{RGB}{10,60,120}
\definecolor{mtgold}{RGB}{180,140,0}
\definecolor{mtgray}{RGB}{90,90,90}
\definecolor{codebg}{RGB}{245,247,250}
\definecolor{codekw}{RGB}{0,80,160}
\definecolor{codestr}{RGB}{0,120,60}
\definecolor{codecom}{RGB}{120,120,120}

% ─── Hyperref ────────────────────────────────────────────────────────────────
\hypersetup{
  colorlinks=true,
  linkcolor=mtblue,
  citecolor=mtblue,
  urlcolor=mtblue,
  pdftitle={Multiplicity Theory for Mathematical Biology: A Comprehensive Framework and Defensive Publication},
  pdfauthor={Multiplicity Foundation},
}

% ─── Listings ────────────────────────────────────────────────────────────────
\lstset{
  backgroundcolor=\color{codebg},
  basicstyle=\ttfamily\small,
  keywordstyle=\color{codekw}\bfseries,
  stringstyle=\color{codestr},
  commentstyle=\color{codecom}\itshape,
  numbers=left,
  numberstyle=\tiny\color{mtgray},
  numbersep=8pt,
  breaklines=true,
  frame=single,
  framerule=0.4pt,
  rulecolor=\color{mtgray!40},
  tabsize=4,
  showstringspaces=false,
  captionpos=b
}

% ─── Theorem environments ─────────────────────────────────────────────────────
\newtheorem{definition}{Definition}[section]
\newtheorem{theorem}{Theorem}[section]
\newtheorem{proposition}{Proposition}[section]
\newtheorem{corollary}{Corollary}[section]
\newtheorem{remark}{Remark}[section]

% ─── Defensive Publication box ───────────────────────────────────────────────
\newtcolorbox{disclosurebox}{
  colback=mtblue!4,
  colframe=mtblue!60,
  fonttitle=\bfseries\color{mtblue},
  title={Prior Art Disclosure Notice},
  breakable
}

% ─── Section formatting ───────────────────────────────────────────────────────
\titleformat{\section}{\large\bfseries\color{mtblue}}{{\thesection}}{1em}{}
\titleformat{\subsection}{\normalsize\bfseries\color{mtgray}}{\thesubsection}{1em}{}
\titleformat{\subsubsection}{\normalsize\itshape\color{mtgray}}{\thesubsubsection}{1em}{}

% ─── Header / Footer ─────────────────────────────────────────────────────────
\pagestyle{fancy}
\fancyhf{}
\fancyhead[L]{\small\color{mtgray}Multiplicity Theory for Mathematical Biology}
\fancyhead[R]{\small\color{mtgray}Defensive Publication --- April 2026}
\fancyfoot[C]{\small\thepage}
\renewcommand{\headrulewidth}{0.4pt}

% ─── Macros ───────────────────────────────────────────────────────────────────
\newcommand{\MT}{\textsc{Multiplicity Theory}}
\newcommand{\Mcal}{\mathcal{M}}
\newcommand{\Pcal}{\mathcal{P}}
\newcommand{\bp}{\mathbf{p}}
\newcommand{\bq}{\mathbf{q}}
\newcommand{\psi}{\psi}
\newcommand{\Psi}{\Psi}
\newcommand{\Phi}{\Phi}
\newcommand{\lam}{\lambda}

% ─────────────────────────────────────────────────────────────────────────────
\begin{document}

% ─── Title page ──────────────────────────────────────────────────────────────
\begin{titlepage}
\centering
\vspace*{1.5cm}

{\color{mtblue}\rule{\linewidth}{1.5pt}}\\[0.4em]
{\color{mtblue}\rule{\linewidth}{0.4pt}}\\[1.5em]

{\Huge\bfseries\color{mtblue}
Multiplicity Theory for\\[0.3em]Mathematical Biology}\\[1em]

{\Large\color{mtgray}
A Prime-Indexed, Recursively Stable Framework\\
for Multi-Scale Biological Modeling}\\[2em]

{\large\color{mtblue}\rule{0.6\linewidth}{0.4pt}}\\[1em]

{\large
\textbf{Comprehensive Report \&}\\
\textbf{Defensive Prior Art Publication}\\[1.5em]
\textcolor{mtgray}{Multiplicity Foundation}\\
\textcolor{mtgray}{Bridgeville, Pennsylvania, USA}\\[1em]
\textcolor{mtgray}{\textit{Publication Date: April 25, 2026}}\\
\textcolor{mtgray}{\textit{Document ID: MTF-BIOL-2026-001}}
}\\[2em]

{\color{mtblue}\rule{0.6\linewidth}{0.4pt}}\\[2em]

\begin{tcolorbox}[colback=mtblue!6,colframe=mtblue!50,width=0.85\linewidth]
\centering\small
\textbf{Defensive Publication Notice.}\\
This document is intentionally published to establish prior art as of its
publication date. All concepts, methods, equations, and frameworks described
herein are placed in the public domain. No patent rights are asserted on the
technical content disclosed. Any subsequent filings by third parties on
substantially identical subject matter are hereby pre-empted by this disclosure.
\end{tcolorbox}

\vfill
{\color{mtblue}\rule{\linewidth}{0.4pt}}\\[0.3em]
{\color{mtblue}\rule{\linewidth}{1.5pt}}
\end{titlepage}

% ─── Abstract ─────────────────────────────────────────────────────────────────
\begin{abstract}
\setlength{\parindent}{0pt}
\onehalfspacing
\noindent
We present \textit{Multiplicity Theory for Mathematical Biology} (MT-MB): a
prime-indexed, recursively stable framework that re-casts classical mathematical
biology---population dynamics, epidemiology, neuroscience, ecological networks,
evolutionary dynamics, reaction-diffusion systems, systems biology, and clinical
laboratory quality control---within a single, unified modeling language.

The core innovation is the replacement of generic tensor indices with
\emph{prime-labeled modes} whose multiplicative structure (via prime
factorization) encodes cross-scale, cross-domain biological configurations.
Biological states are represented as tensors over a multiplicity space
$\Mcal = \bigoplus_{\bp} \mathbb{R}$, governed by recursive updates coupled with
explicit \emph{prime-wise aggregate stability constraints}. A global
meta-equation $H(t,\psi(t)) \to M(t,\psi(t))\,T(t,G) + f(t,\psi(t)) =
\lambda(t)\psi(t)$ unifies all sectors and identifies four interpretable
components: multiplicity constraints ($M$), structural networks ($T$), exogenous
drivers ($f$), and global consistency ($\lambda$).

The framework generates three families of testable predictions across all
domains: structured correlations along prime modes, emergent interaction
motifs, and early-warning signals for regime transitions. Python reference
implementations for three key modules (clinical lab quality control,
epidemiology, and population dynamics) are provided. This document simultaneously
serves as a comprehensive technical report and a \textbf{defensive prior art
publication}, placing the full intellectual content in the public domain.
\end{abstract}

\vspace{1em}
\noindent\textbf{Keywords:} Multiplicity Theory; prime indexing; tensor
dynamics; recursive stability; mathematical biology; SIR models; Lotka--Volterra;
neural networks; systems biology; reaction-diffusion; evolutionary dynamics;
clinical laboratory systems; defensive publication; prior art.
\newpage
\tableofcontents

% ═══════════════════════════════════════════════════════════════════════════════
\section{Executive Summary}
% ═══════════════════════════════════════════════════════════════════════════════

\MT{} (MT) is a mathematical framework in which the fundamental objects of
inquiry are \emph{multiplicities}: prime-indexed, recursively generated patterns
of interaction. Rather than treating mathematical objects as static entities, MT
reinterprets them as recursively generated patterns of prime-labeled interactions,
where sets and modules become \emph{multiplicity spaces} whose identity and
behavior are preserved through recursive feedback across scales.

\subsection*{What This Document Establishes}

This report synthesizes and formalizes the application of \MT{} to the eight
core areas of mathematical biology:

\begin{enumerate}[leftmargin=*]
  \item \textbf{Population Dynamics} --- Lotka--Volterra equations promoted to
    prime-indexed population tensors $P_{\bp}(t)$ with species, age, genotype,
    niche, and resource primes.
  \item \textbf{Epidemiological Models} --- SIR and its variants reformulated as
    flows among prime-labeled compartments (host, pathogen, setting, testing).
  \item \textbf{Neural Networks and Recursive Systems} --- Hodgkin--Huxley and
    rate-based dynamics on prime-labeled (neuron, dendrite, circuit, region,
    neuromodulator, synapse) configurations with quantum-inspired superposition
    on a multiplicity basis.
  \item \textbf{Ecological Networks} --- Food webs and resource-competition
    modeled via multi-layer prime-indexed energy-flow tensors.
  \item \textbf{Evolutionary Dynamics} --- Wright--Fisher and replicator
    equations on prime-labeled (locus, allele, background, environment, deme)
    tensors, accommodating mutation, selection, recombination, and drift.
  \item \textbf{Reaction--Diffusion and Turing Patterns} --- Spatiotemporal
    pattern formation via prime-indexed concentration tensors with
    diffusion encoded as coupling between spatial primes.
  \item \textbf{Systems Biology} --- Metabolite, gene expression, and enzyme
    activity tensors over shared prime-labeled configurations, enabling
    cross-scale metabolic/regulatory/activity feedback.
  \item \textbf{Clinical Laboratory Systems} --- Error-propagation and quality
    control modeled via prime modes for machine, reagent, operator, and patient
    class with recursive stability targeting lab-acceptable baselines.
\end{enumerate}

\subsection*{Key Innovations}

\begin{itemize}[leftmargin=*]
  \item \textbf{Prime-factorized index sets}: Replacing generic tensor indices
    with multiplicative prime labels turns the index set into a monoid, enabling
    algebraic splitting, merging, and projection operations that map onto real
    biological processes (hybridization, cell division, cross-sector interaction).
  \item \textbf{Recursive stability constraints}: Adding prime-wise balance
    equations $\mathcal{X}_p(t+1) - \mathcal{X}_p(t) = \Phi_p(X(t))$ to bare
    recursions yields constrained dynamical systems where changes in one
    configuration must be compensated by others sharing the same prime,
    enforcing stable long-term behavior at every scale.
  \item \textbf{Unified meta-equation}: A single operator equation
    $H \to MT + f = \lambda\psi$ governs all sectors, with $M$, $T$, $f$,
    $\lambda$ admitting concrete domain-specific interpretations.
  \item \textbf{Grounded quantum-inspired representation}: The quantum analogy
    $\Psi(t) = \sum_{\bp} a_{\bp}(t)\,\psi_{\bp}$ is given an explicit
    multiplicity basis, turning it from a metaphor into a genuine Hilbert space
    spanned by prime-labeled biological configurations.
\end{itemize}

\subsection*{Defensive Publication Purpose}

This document is published \emph{in full} to establish unambiguous prior art
as of April 25, 2026. All mathematical formulations, methodological structures,
algorithmic implementations, and predictive claims documented herein are placed
in the public domain under the principles of defensive disclosure. Third-party
patent applications on substantially identical subject matter filed after this
date are pre-empted.

\newpage
% ═══════════════════════════════════════════════════════════════════════════════
\section{Multiplicity Framework for Mathematical Biology}
\label{sec:framework}
% ═══════════════════════════════════════════════════════════════════════════════

\subsection{Prime-Labeled Multiplicity Space}

\begin{definition}[Prime mode set]
Let $\Pcal = \{p_1, p_2, \ldots\}$ be a set of distinct prime numbers, each
assigned to a biologically meaningful factor: species, host type, brain region,
lab machine, gene locus, spatial patch, etc.
\end{definition}

\begin{definition}[Configuration]
A \emph{configuration} is a finite multiset of primes:
\[
  \bp = \bigl(p_1^{e_1},\, p_2^{e_2},\, \ldots,\, p_k^{e_k}\bigr),
  \quad e_i \in \mathbb{Z}_{\geq 0},
\]
with integer index $n(\bp) = \prod_i p_i^{e_i}$ as its canonical label.
The set of all configurations is the free commutative monoid $\Pcal^*$ over
$\Pcal$.
\end{definition}

\begin{definition}[Multiplicity space]
\[
  \Mcal \;=\; \bigoplus_{\bp \in \Pcal^*} \mathbb{R}.
\]
Each coordinate axis of $\Mcal$ corresponds to one configuration $\bp$.
Algebraic operations on $\Pcal^*$ (multiplication, factorization) induce
natural maps on $\Mcal$ that model biological processes such as hybridization,
cell division, or cross-sector interaction.
\end{definition}

\subsection{State Tensors and Recursive Updates}

A biological system is described by a \emph{state tensor}
\[
  X_{\bp}(t) \in \mathbb{R}, \quad \bp \in \Pcal^*,
\]
representing population size, infection level, neural activity, concentration,
error magnitude, or any domain-specific quantity for configuration $\bp$ at
time $t$.

\paragraph{Discrete-time dynamics.}
\begin{equation}
  X_{\bp}(t+1) \;=\; X_{\bp}(t) \;+\; \Delta_{\bp}\!\bigl(X(t),\,\theta(t)\bigr),
  \label{eq:discrete}
\end{equation}
where $\Delta_{\bp}$ encodes creation, removal, interaction, and movement
in configuration $\bp$; $\theta(t)$ are domain-specific parameters.

\paragraph{Continuous-time dynamics.}
\begin{equation}
  \frac{\partial X_{\bp}}{\partial t}
  \;=\; D_{\bp}\,\nabla^2 X_{\bp} \;+\; F_{\bp}(X(t)),
  \label{eq:continuous}
\end{equation}
with $D_{\bp}$ an effective diffusion coefficient and $F_{\bp}$ local
reaction kinetics (applicable to reaction-diffusion sectors).

Interaction structures are themselves tensors over pairs or tuples of
configurations, e.g., $K_{\bp,\bq}$ for pairwise coupling from $\bq$ to
$\bp$.

\subsection{Prime-Wise Aggregates and Recursive Stability}

\begin{definition}[Prime aggregate]
For each prime mode $p \in \Pcal$:
\begin{equation}
  \mathcal{X}_p(t) \;=\; \sum_{\bp \ni p} X_{\bp}(t).
  \label{eq:aggregate}
\end{equation}
\end{definition}

\begin{definition}[Recursive stability constraint]
\begin{equation}
  \mathcal{X}_p(t+1) - \mathcal{X}_p(t)
  \;=\; \Phi_p\!\bigl(X(t)\bigr),
  \label{eq:stability}
\end{equation}
where $\Phi_p : \Mcal \to \mathbb{R}$ is chosen so that $\mathcal{X}_p(t)$
converges to a prescribed long-term target (fixed point, limit cycle, or
other attractor).
\end{definition}

\begin{remark}
Equation~(\ref{eq:stability}) turns the bare recursion~(\ref{eq:discrete}) into
a \emph{constrained dynamical system}: any change in one configuration is
balanced by compensatory changes in all other configurations sharing the same
prime mode, preserving prime-wise stability globally.
\end{remark}

\subsection{Global Multiplicity State and Meta-Equation}

All sector-specific tensors are assembled into a global state:
\[
  \psi(t) \;\in\; \Mcal_{\mathrm{global}}
  \;=\; \bigoplus_{\text{sectors}} \Mcal_{\text{sector}}.
\]

\begin{definition}[Multiplicity meta-equation]
\begin{equation}
  H\bigl(t,\psi(t)\bigr)
  \;\longrightarrow\;
  M\bigl(t,\psi(t)\bigr)\,T(t,G)
  \;+\; f\bigl(t,\psi(t)\bigr)
  \;=\; \lambda(t)\,\psi(t),
  \label{eq:meta}
\end{equation}
where:
\begin{itemize}[leftmargin=*]
  \item $M(t,\psi(t))$: operator encoding all prime-wise stability constraints
    $\{\Phi_p\}$ across sectors (multiplicity constraints),
  \item $T(t,G)$: structural interaction networks (ecological tensors, contact
    networks, synaptic connectivity, reaction stoichiometry, lab workflows)
    parameterized by network graph $G$,
  \item $f(t,\psi(t))$: exogenous drivers (seasonality, interventions,
    external stimuli, noise),
  \item $\lambda(t)$: global consistency multiplier enforcing conservation,
    scaling laws, or cross-sector coupling.
\end{itemize}
\end{definition}

\subsection{Quantum-Inspired Multiplicity Superposition}

For domains where superposition or uncertainty is central (neural dynamics,
genetic evolution, cellular states):

\begin{equation}
  \Psi(t)
  \;=\; \sum_{\bp \in \Pcal^*} a_{\bp}(t)\,\psi_{\bp},
  \label{eq:superposition}
\end{equation}
where $\{\psi_{\bp}\}$ is the orthonormal basis of prime-labeled configurations
and $|a_{\bp}(t)|^2$ gives the probability or intensity of configuration $\bp$
at time $t$. The recursive update~(\ref{eq:discrete}) and
constraint~(\ref{eq:stability}) induce well-defined dynamics on $a_{\bp}(t)$.

\subsection{Generic Predictions and Validation Template}

Across all domains, the MT-MB framework produces three families of predictions:

\begin{enumerate}[leftmargin=*]
  \item \textbf{Structured correlations along prime modes}: observables covary
    within prime-sharing configuration sets, revealing low-dimensional
    ``prime channels'' in high-dimensional data.
  \item \textbf{Emergent interaction motifs}: configuration-specific interaction
    tensors expose strong cross-scale couplings invisible to mean-field models.
  \item \textbf{Early-warning signals for regime transitions}: prime aggregates
    $\mathcal{X}_p(t)$ approaching stability-envelope boundaries anticipate
    regime shifts (epidemics, tipping points, seizures, lab failures) before
    observable changes in individual variables.
\end{enumerate}

The shared validation protocol is:
\begin{enumerate}[leftmargin=*]
  \item Select a domain and dataset with known multi-factor structure.
  \item Choose 2--3 prime modes corresponding to empirically observable factors.
  \item Fit the prime-constrained recursive multiplicity model.
  \item Compare predictive performance against domain-standard baselines.
\end{enumerate}

\newpage
% ═══════════════════════════════════════════════════════════════════════════════
\section{Domain Instantiations}
\label{sec:domains}
% ═══════════════════════════════════════════════════════════════════════════════

\subsection{Population Dynamics}

\paragraph{Prime modes.}
$p_{\text{s}}$ (species), $p_{\text{a}}$ (age/life stage), $p_{\text{g}}$
(genotype/phenotype), $p_{\text{n}}$ (niche/patch), $p_{\text{r}}$ (resource).

\paragraph{Population tensor and update.}
\[
  P_{\bp}(t+1)
  = P_{\bp}(t)
  + \Delta_{\bp}\!\bigl(P(t),\,\alpha,\,\beta,\,\gamma\bigr),
\]
with $\Delta_{\bp}$ encoding births, deaths, predation, competition, and
dispersal. Predator-prey interactions become:
\[
  \frac{\partial P_i}{\partial t}
  = \alpha_i P_i
  - \sum_{\bq} \beta_{i\bq}\,P_i\,P_{\bq},
\]
where $\beta_{i\bq}$ is the predation coefficient tensor.

\paragraph{Prime aggregate and stability.}
\[
  \mathcal{P}_p(t) = \sum_{\bp \ni p} P_{\bp}(t),
  \qquad
  \mathcal{P}_p(t+1) - \mathcal{P}_p(t) = \Phi_p^{(\mathrm{pop})}(P(t)).
\]
$\Phi_p$ drives total species biomass toward carrying capacity or a
predator-prey limit cycle.

\paragraph{Meta-equation mapping.}
$T$ = ecological interaction tensor; $f$ = climate/resource pulses;
$\lambda$ = biomass/energy scaling law.

\subsection{Epidemiological Models}

\paragraph{Prime modes.}
$p_{\text{h}}$ (host type), $p_{\text{p}}$ (pathogen strain), $p_{\text{s}}$
(contact setting), $p_{\text{t}}$ (testing/reporting modality).

\paragraph{Infection tensor and update.}
\[
  I_{\bp}(t+1)
  = I_{\bp}(t)
  + \Delta_{\bp}\!\bigl(I(t),\,S(t),\,R(t),\,\beta,\,\gamma\bigr),
\]
with $\beta_{\bp}$ a configuration-dependent transmission rate. The classical
SIR system,
\[
  \frac{dS}{dt} = -\beta SI,\quad
  \frac{dI}{dt} = \beta SI - \gamma I,\quad
  \frac{dR}{dt} = \gamma I,
\]
is recovered as the marginal of the multiplicity dynamics projected onto
configurations containing only $p_{\text{h}}$ (a single, homogeneous host
type).

\paragraph{Prime aggregate and stability.}
\[
  \mathcal{I}_p(t) = \sum_{\bp \ni p} I_{\bp}(t),
  \qquad
  \mathcal{I}_p(t+1) - \mathcal{I}_p(t) = \Phi_p^{(\mathrm{epi})}(I,S,R).
\]
$\Phi_p$ targets an endemic equilibrium, a seasonal limit cycle, or
elimination.

\paragraph{Meta-equation mapping.}
$T$ = contact network $\beta_{\bp,\bq}$; $f$ = seasonality, interventions;
$\lambda$ = population-size conservation.

\subsection{Neural Networks and Recursive Systems}

\paragraph{Prime modes.}
$p_{\text{n}}$ (neuron), $p_{\text{d}}$ (dendrite), $p_{\text{c}}$ (circuit),
$p_{\text{r}}$ (brain region), $p_{\text{m}}$ (neuromodulator), $p_{\text{s}}$
(synapse type).

\paragraph{Activity tensor and update.}
\[
  N_{\bp}(t+1)
  = N_{\bp}(t)
  + \Delta_{\bp}\!\bigl(N(t),\,W(t),\,\theta\bigr),
\]
with $W_{\bp,\bq}(t)$ the synaptic weight tensor and $\Delta_{\bp}$ encoding
Hodgkin--Huxley or integrate-and-fire dynamics in configuration $\bp$.

\paragraph{Quantum-inspired superposition.}
\[
  \Psi_N(t) = \sum_{\bp} a_{\bp}(t)\,\psi_{\bp},
  \quad
  |a_{\bp}(t)|^2 \;=\; \text{activation intensity of configuration }\bp.
\]

\paragraph{Meta-equation mapping.}
$T$ = synaptic connectivity $W_{\bp,\bq}$; $f$ = sensory stimuli,
neuromodulation; $\lambda$ = metabolic/energy constraint.

\subsection{Ecological Networks}

\paragraph{Prime modes.}
$p_{\text{s}}$ (species), $p_{\text{t}}$ (trophic level), $p_{\text{h}}$
(habitat type), $p_{\text{r}}$ (resource class).

\paragraph{Energy-flow tensor.}
\[
  E_{\bp}(t) = \text{energy/biomass flow at configuration }\bp,
\]
with harmonic cycles emerging from recursive feedback:
\[
  E(t) = A\sin\!\left(\tfrac{2\pi t}{T}\right) + B\cos\!\left(\tfrac{2\pi t}{T}\right),
\]
where $T$ is the ecological period. Resource dynamics obey:
\[
  R^{(t+1)} = R^{(t)} - \Delta R + f(E_{\bp}).
\]

\paragraph{Meta-equation mapping.}
$T$ = food-web adjacency tensor; $f$ = environmental pulses; $\lambda$ =
energy-conservation constraint.

\subsection{Evolutionary Dynamics}

\paragraph{Prime modes.}
$p_{\text{l}}$ (locus), $p_{\text{a}}$ (allele), $p_{\text{b}}$ (epistatic
background), $p_{\text{e}}$ (environment), $p_{\text{d}}$ (demographic class).

\paragraph{Genetic tensor and update.}
\[
  G_{\bp}(t+1)
  = G_{\bp}(t)
  + \Delta_{\bp}\!\bigl(G(t),\,\mu(t),\,s(t),\,D(t)\bigr),
\]
with mutation tensor $M_{\bp\to\bq}$, fitness tensor $W_{\bp}(t)$, and
stochastic drift modeled as noise on $G_{\bp}(t)$.

\paragraph{Prime aggregate and stability.}
\[
  \mathcal{G}_p(t) = \sum_{\bp \ni p} G_{\bp}(t),
  \qquad
  \mathcal{G}_p(t+1) - \mathcal{G}_p(t) = \Phi_p^{(\mathrm{evo})}(G(t)).
\]
$\Phi_p$ encodes mutation-selection balance, balanced polymorphism, or Red
Queen dynamics.

\paragraph{Quantum-inspired genetic superposition.}
\[
  \Psi_G(t) = \sum_{\bp} c_{\bp}(t)\,\psi_{\bp},
  \quad
  |c_{\bp}(t)|^2 \approx G_{\bp}(t)/N \text{ (large-}N\text{ limit)}.
\]

\paragraph{Meta-equation mapping.}
$T$ = mutation/recombination/fitness tensors; $f$ = environmental shifts,
bottlenecks; $\lambda$ = population-size normalization.

\subsection{Reaction--Diffusion and Turing Patterns}

\paragraph{Prime modes.}
$p_{\text{s}}$ (chemical species), $p_{\text{c}}$ (compartment),
$p_{\text{d}}$ (domain/tissue type), $p_{\text{x}}$ (spatial patch).

\paragraph{Concentration tensor and PDE.}
\begin{equation}
  \frac{\partial C_{\bp}}{\partial t}
  = D_{\bp}\,\nabla^2 C_{\bp} + F_{\bp}(C(t)),
  \label{eq:rd}
\end{equation}
where on a discrete lattice $\nabla^2$ becomes coupling between configurations
differing only in $p_{\text{x}}$. Reaction kinetics:
\[
  F_{\bp}(C) = \sum_{\bq_1,\ldots,\bq_k}
    R_{\bp,\bq_1,\ldots,\bq_k}\prod_{j=1}^k C_{\bq_j}(t).
\]

\paragraph{Prime aggregate (continuous).}
\[
  \mathcal{C}_p(t) = \int_\Omega \sum_{\bp \ni p} C_{\bp}(x,t)\,dx,
  \qquad
  \frac{d}{dt}\,\mathcal{C}_p = \Phi_p^{(\mathrm{rd})}(C(t)).
\]

\paragraph{Modified Turing analysis.} The homogeneous steady state is the
configuration where all spatial primes share equal $C_{\bp}$. Prime-wise
constraints restrict which spatial instabilities are admissible, sharpening
classical Turing conditions.

\paragraph{Meta-equation mapping.}
$T$ = reaction tensors $R$ plus diffusion coupling $D_{\bp,\bq}$;
$f$ = morphogen sources, boundary conditions; $\lambda$ = mass conservation.

\subsection{Systems Biology}

\paragraph{Prime modes.}
$p_{\text{m}}$ (metabolite), $p_{\text{r}}$ (reaction/enzyme),
$p_{\text{g}}$ (gene/module), $p_{\text{c}}$ (compartment),
$p_{\text{t}}$ (cell type), $p_{\text{s}}$ (spatial/tissue region).

\paragraph{State tensors.}
\[
  M_{\bp}(t),\quad E_{\bp}(t),\quad A_{\bp}(t),
\]
for metabolite concentration, gene expression, and enzyme activity,
respectively, over $\Mcal = \bigoplus_{\bp} \mathbb{R}^3$.

\paragraph{Coupled recursive updates.}
\begin{align*}
  M_{\bp}(t+1) &= M_{\bp}(t) + \Delta_{\bp}^{(\mathrm{metab})}(M,E,A,S),\\
  E_{\bp}(t+1) &= E_{\bp}(t) + \Delta_{\bp}^{(\mathrm{reg})}(E,M,\Theta),\\
  A_{\bp}(t+1) &= A_{\bp}(t) + \Delta_{\bp}^{(\mathrm{act})}(A,E,M),
\end{align*}
with stoichiometry tensor $S_{\bp,\bq}$ and regulatory tensor
$\Theta_{\bp,\bq}$.

\paragraph{Prime aggregates and stability.}
\begin{align*}
  \mathcal{M}_p(t+1) - \mathcal{M}_p(t)
    &= \Phi_p^{(\mathrm{metab})}(M,E,A),\\
  \mathcal{E}_p(t+1) - \mathcal{E}_p(t)
    &= \Phi_p^{(\mathrm{reg})}(E,M),\\
  \mathcal{A}_p(t+1) - \mathcal{A}_p(t)
    &= \Phi_p^{(\mathrm{act})}(A,E,M).
\end{align*}

\paragraph{Meta-equation mapping.}
$T$ = reaction stoichiometry $S_{\bp,\bq}$ and regulatory network
$\Theta_{\bp,\bq}$; $f$ = nutrient availability, drug treatments;
$\lambda$ = mass/energy and thermodynamic consistency.

\subsection{Clinical Laboratory Systems}

\paragraph{Prime modes.}
$p_{\text{m}}$ (machine/analyzer), $p_{\text{r}}$ (reagent batch),
$p_{\text{o}}$ (operator/technician), $p_{\text{c}}$ (patient class).

\paragraph{Error tensor and update.}
\[
  E_{\bp}(t+1) = E_{\bp}(t)
  + \Delta_{\bp}\!\bigl(E(t),\,\mathrm{calib}(t)\bigr),
\]
where $\Delta_{\bp}$ captures drift, recalibration events, and systematic
machine--reagent bias.

\paragraph{Prime aggregate and stability.}
\[
  \mathcal{E}_p(t) = \sum_{\bp \ni p} E_{\bp}(t),
  \qquad
  \mathcal{E}_p(t+1) - \mathcal{E}_p(t) = \Phi_p^{(\mathrm{lab})}(E(t)).
\]
$\Phi_p$ drives error toward the laboratory's acceptable baseline; crossing
a stability boundary signals an out-of-control event.

\paragraph{Meta-equation mapping.}
$T$ = lab workflow and equipment dependency graph; $f$ = reagent lot changes,
staff turnover; $\lambda$ = cross-mode consistency (total error budget).

\newpage
% ═══════════════════════════════════════════════════════════════════════════════
\section{Comprehensive Mathematical Overview}
\label{sec:math}
% ═══════════════════════════════════════════════════════════════════════════════

\subsection{Algebraic Structure of the Multiplicity Index Set}

The configuration set $\Pcal^* = \{\bp = \prod_i p_i^{e_i} : e_i \geq 0\}$
is the free commutative monoid generated by $\Pcal$. The natural number
$n(\bp) = \prod_i p_i^{e_i}$ provides a unique integer label by the
Fundamental Theorem of Arithmetic.

\begin{proposition}
The map $n : \Pcal^* \to \mathbb{Z}_{>0}$ is a monoid isomorphism onto the
multiplicative monoid of positive integers. In particular, the
\emph{divisibility partial order} on $\mathbb{Z}_{>0}$ induces a
\emph{containment partial order} on configurations:
\[
  \bp \leq \bq \;\iff\; n(\bp) \mid n(\bq) \;\iff\;
  e_i(\bp) \leq e_i(\bq)\;\forall\,i.
\]
\end{proposition}

This means ``$\bp$ is a sub-configuration of $\bq$'' is precisely divisibility,
and prime projections (e.g., setting $e_i = 0$) correspond to marginalizing
over mode $p_i$.

\subsection{Multiplicity Space as a Graded Module}

\begin{definition}
The multiplicity space $\Mcal = \bigoplus_{\bp} \mathbb{R}$ carries a natural
grading by $n(\bp)$:
\[
  \Mcal = \bigoplus_{n=1}^\infty \Mcal_n,
  \quad
  \Mcal_n = \bigoplus_{n(\bp)=n} \mathbb{R}.
\]
The grade-$1$ component $\Mcal_1 \cong \mathbb{R}$ is the trivial configuration
(no modes active). Grade-$p$ for prime $p$ gives the single-mode states.
\end{definition}

\begin{remark}
Tensor products of multiplicity spaces correspond to joining biological
systems: $\Mcal_{\text{species}} \otimes \Mcal_{\text{space}}$ yields the
combined species-spatial state space. Prime factorization then provides
automatic decomposition of cross-system dynamics.
\end{remark}

\subsection{Fixed-Point and Stability Theory}

Let $\Phi = (\Phi_p)_{p \in \Pcal}$ be the collection of stability functions.
Define the \emph{aggregate flow}:
\[
  \Phi : \Mcal \to \mathbb{R}^{|\Pcal|},
  \quad
  \Phi(X)_p = \Phi_p(X).
\]

\begin{definition}[Recursive stability]
A state $X^* \in \Mcal$ is \emph{recursively stable} if
$\Phi_p(X^*) = 0$ for all $p \in \Pcal$, and the Jacobian
$D\Phi|_{X^*}$ has all eigenvalues with negative real part.
\end{definition}

\begin{theorem}[Existence of recursively stable states]
Under Lipschitz continuity of $\Phi_p$ and the sector-specific $\Delta_{\bp}$,
and assuming bounded interaction tensors $K_{\bp,\bq}$, the dynamical system
\eqref{eq:discrete}--\eqref{eq:stability} admits at least one fixed point
$X^* \in \Mcal$.
\end{theorem}

\begin{corollary}[Prime-channel decomposition]
At a recursively stable state $X^*$, the linearized dynamics decouple into
independent prime-channel modes. Perturbations in mode $\bp$ propagate only
to modes $\bq$ with $\gcd(n(\bp), n(\bq)) > 1$ (shared prime factors),
predicting the \emph{prime channel structure} of observable covariance.
\end{corollary}

\subsection{Spectral Theory for Prime Aggregates}

Let $\mathbf{A}$ be the aggregate matrix with entries
$A_{p,\bp} = \mathbf{1}[p \mid n(\bp)]$. The aggregate vector is:
\[
  \boldsymbol{\mathcal{X}}(t)
  = \mathbf{A}\,\mathbf{X}(t)
  \in \mathbb{R}^{|\Pcal|},
\]
where $\mathbf{X}(t) \in \mathbb{R}^{|\Pcal^*|}$ is the full state vector.

The stability constraint becomes:
\[
  \boldsymbol{\mathcal{X}}(t+1) - \boldsymbol{\mathcal{X}}(t)
  = \boldsymbol{\Phi}(\mathbf{X}(t)).
\]
For linear $\boldsymbol{\Phi}(\mathbf{X}) = \mathbf{B}\,\mathbf{X}$ with
stability matrix $\mathbf{B}$, the system is globally stable iff all
eigenvalues of $\mathbf{A}^+\mathbf{B}$ have negative real part, where
$\mathbf{A}^+$ is the Moore-Penrose pseudoinverse of $\mathbf{A}$.

\subsection{Connection to the Meta-Equation}

Equation~\eqref{eq:meta} can be written in operator form. Let $\mathcal{H}$
be the Hilbert space with basis $\{\psi_{\bp}\}$. Define:

\begin{itemize}[leftmargin=*]
  \item \textbf{Multiplicity operator} $\hat{M}[\psi] = \sum_p \Phi_p(\psi)\,
    \partial/\partial\mathcal{X}_p$: encodes how stability constraints
    constrain the flow of $\psi(t)$.
  \item \textbf{Transfer operator} $\hat{T}[G]\,\psi = \sum_{\bp,\bq}
    K_{\bp,\bq}\,\psi_{\bq}$: acts on states via interaction tensors.
  \item \textbf{Forcing operator} $\hat{f}[t,\psi] = f(t)\,\psi + \xi(t)$:
    adds exogenous perturbations and noise.
\end{itemize}

The meta-equation is then the operator eigenvalue problem:
\[
  \bigl(\hat{M} + \hat{T}\bigr)\psi + \hat{f} = \lambda(t)\,\psi,
\]
whose eigenvectors $\psi$ describe the stable configurations of the whole
system, and whose eigenvalues $\lambda(t)$ encode the global consistency or
``fitness'' of the current state.

\subsection{Cross-Domain Coupling}

When multiple biological sectors (e.g., epidemiology and population dynamics)
share primes (e.g., $p_{\text{host type}}$), their state tensors are coupled:

\[
  \psi(t) = \psi_{\text{epi}}(t) \oplus \psi_{\text{pop}}(t),
  \quad
  \mathcal{X}_p^{\text{total}}(t)
  = \mathcal{I}_p(t) + \mathcal{P}_p(t).
\]

Cross-domain stability constraints then link epidemic dynamics to population
structure, enabling joint modeling of eco-epidemiological systems without
ad hoc coupling terms.

\newpage
% ═══════════════════════════════════════════════════════════════════════════════
\section{Reference Code Implementations}
\label{sec:code}
% ═══════════════════════════════════════════════════════════════════════════════

The following Python code provides reference implementations of three key
modules. They are provided as proof of concept and to fix the computational
interpretation of the mathematical definitions. All implementations are
released into the public domain.

\subsection{Core Multiplicity Space}

\begin{lstlisting}[language=Python, caption={Core multiplicity space and prime-mode infrastructure.}]
"""
multiplicity_core.py
Core data structures for the Multiplicity Theory framework for
Mathematical Biology.
Released into the public domain -- Multiplicity Foundation, 2026.
"""

from __future__ import annotations
from dataclasses import dataclass, field
from typing import Callable, Dict, List, Set, Tuple
from math import prod
import numpy as np
from sympy import factorint


# ---------------------------------------------------------------------------
# Configuration (prime multiset)
# ---------------------------------------------------------------------------

@dataclass(frozen=True)
class Configuration:
    """
    A biological configuration: a finite multiset of prime modes.
    Internally stored as a frozenset of (prime, exponent) pairs.
    The canonical integer label is prod(p**e for p,e in factors).
    """
    factors: frozenset  # frozenset of (prime: int, exponent: int)

    @classmethod
    def from_dict(cls, d: Dict[int, int]) -> "Configuration":
        """Create from {prime: exponent} dict."""
        return cls(frozenset((p, e) for p, e in d.items() if e > 0))

    @classmethod
    def from_int(cls, n: int) -> "Configuration":
        """Create from canonical integer label (prime factorization)."""
        return cls.from_dict(factorint(n))

    @property
    def label(self) -> int:
        return prod(p**e for p, e in self.factors)

    def contains_prime(self, p: int) -> bool:
        return any(prime == p for prime, _ in self.factors)

    def primes(self) -> Set[int]:
        return {p for p, _ in self.factors}

    def __mul__(self, other: "Configuration") -> "Configuration":
        """Join two configurations (multiply labels)."""
        d: Dict[int, int] = {}
        for p, e in self.factors:
            d[p] = d.get(p, 0) + e
        for p, e in other.factors:
            d[p] = d.get(p, 0) + e
        return Configuration.from_dict(d)

    def __repr__(self) -> str:
        parts = sorted(f"p{p}^{e}" for p, e in self.factors)
        return "(" + " * ".join(parts) + ")" if parts else "(empty)"


# ---------------------------------------------------------------------------
# Multiplicity Space
# ---------------------------------------------------------------------------

class MultiplicitySpace:
    """
    The multiplicity space M = direct_sum_{p in P*} R.
    Implemented as a dictionary {Configuration -> float}.
    """

    def __init__(self, prime_set: List[int]):
        self.prime_set = prime_set          # available prime modes
        self.state: Dict[Configuration, float] = {}

    def set(self, config: Configuration, value: float) -> None:
        self.state[config] = value

    def get(self, config: Configuration) -> float:
        return self.state.get(config, 0.0)

    def aggregate(self, prime: int) -> float:
        """Compute X_p(t) = sum_{p in bp} X_bp(t)."""
        return sum(
            v for c, v in self.state.items()
            if c.contains_prime(prime)
        )

    def all_configs(self) -> List[Configuration]:
        return list(self.state.keys())

    def as_array(self, configs: List[Configuration]) -> np.ndarray:
        return np.array([self.get(c) for c in configs])


# ---------------------------------------------------------------------------
# Recursive Stability Constraint
# ---------------------------------------------------------------------------

@dataclass
class StabilityConstraint:
    """
    Encodes Phi_p: M -> R for a single prime mode p.
    The constraint drives aggregate X_p toward target.
    """
    prime: int
    target: float           # desired long-term aggregate value
    gain: float = 0.1       # proportional feedback gain

    def phi(self, space: MultiplicitySpace) -> float:
        current = space.aggregate(self.prime)
        return self.gain * (self.target - current)


# ---------------------------------------------------------------------------
# Multiplicity Dynamical System
# ---------------------------------------------------------------------------

class MultiplicitySystem:
    """
    Implements X_bp(t+1) = X_bp(t) + Delta_bp(X, params)
    subject to prime-wise stability constraints.
    """

    def __init__(
        self,
        space: MultiplicitySpace,
        delta: Callable[[Configuration, MultiplicitySpace, dict], float],
        constraints: List[StabilityConstraint],
        params: dict = None,
    ):
        self.space = space
        self.delta = delta
        self.constraints = constraints
        self.params = params or {}

    def step(self) -> None:
        """Advance one discrete time step."""
        configs = self.space.all_configs()
        # 1. Compute raw updates
        updates = {c: self.delta(c, self.space, self.params) for c in configs}
        # 2. Compute stability corrections
        corrections: Dict[Configuration, float] = {c: 0.0 for c in configs}
        for constraint in self.constraints:
            phi_val = constraint.phi(self.space)
            # Distribute correction equally over configs containing this prime
            relevant = [c for c in configs if c.contains_prime(constraint.prime)]
            if relevant:
                per_config = phi_val / len(relevant)
                for c in relevant:
                    corrections[c] += per_config
        # 3. Apply updates + corrections
        for c in configs:
            old = self.space.get(c)
            self.space.set(c, old + updates[c] + corrections[c])

    def run(self, T: int) -> List[Dict[Configuration, float]]:
        """Run T steps and return list of snapshots."""
        history = []
        for _ in range(T):
            history.append(dict(self.space.state))
            self.step()
        return history
\end{lstlisting}

\subsection{Clinical Laboratory Quality Control Module}

\begin{lstlisting}[language=Python, caption={Prime-encoded clinical lab QC model with recursive stability.}]
"""
lab_qc_multiplicity.py
Prime-encoded clinical laboratory quality control module.
Prime modes: p_m (machine), p_r (reagent), p_o (operator), p_c (patient class)
"""

import numpy as np
from multiplicity_core import (
    Configuration, MultiplicitySpace, StabilityConstraint, MultiplicitySystem
)

# Define prime modes
P_MACHINE   = 2    # machine/analyzer
P_REAGENT   = 3    # reagent batch
P_OPERATOR  = 5    # operator
P_PATIENT   = 7    # patient class

def build_lab_space(error_init: dict) -> MultiplicitySpace:
    """
    error_init: {(machine_id, reagent_id, operator_id, patient_id): error_val}
    where each id is 0 or 1 (present/absent in this first model).
    """
    space = MultiplicitySpace(prime_set=[P_MACHINE, P_REAGENT, P_OPERATOR, P_PATIENT])
    for (m, r, o, c), val in error_init.items():
        factors = {}
        if m: factors[P_MACHINE]   = 1
        if r: factors[P_REAGENT]   = 1
        if o: factors[P_OPERATOR]  = 1
        if c: factors[P_PATIENT]   = 1
        config = Configuration.from_dict(factors)
        space.set(config, val)
    return space


def lab_delta(
    config: Configuration,
    space: MultiplicitySpace,
    params: dict
) -> float:
    """
    Delta_bp for lab error:
    - Drift: Gaussian noise scaled by config-specific drift rate.
    - Cross-mode interaction: machine-reagent bias.
    """
    drift_rate = params.get("drift_rate", 0.02)
    rng = params.get("rng", np.random.default_rng(42))
    drift = rng.normal(0, drift_rate)

    # Machine-reagent systematic bias
    bias = 0.0
    if (config.contains_prime(P_MACHINE) and config.contains_prime(P_REAGENT)):
        bias = params.get("machine_reagent_bias", 0.05)

    return drift + bias


def run_lab_simulation(
    error_init: dict,
    target_errors: dict,  # {prime: target aggregate error}
    T: int = 200,
    seed: int = 0
) -> dict:
    """
    Run T steps of the prime-encoded lab QC model.
    Returns: dict with 'history', 'aggregates', 'early_warnings'.
    """
    rng = np.random.default_rng(seed)
    space = build_lab_space(error_init)

    constraints = [
        StabilityConstraint(prime=p, target=target, gain=0.15)
        for p, target in target_errors.items()
    ]

    params = {
        "drift_rate": 0.03,
        "machine_reagent_bias": 0.04,
        "rng": rng,
    }

    system = MultiplicitySystem(space, lab_delta, constraints, params)
    history = system.run(T)

    # Compute aggregate time series and early warnings
    primes = [P_MACHINE, P_REAGENT, P_OPERATOR, P_PATIENT]
    prime_names = {P_MACHINE: "machine", P_REAGENT: "reagent",
                   P_OPERATOR: "operator", P_PATIENT: "patient"}
    aggregates = {name: [] for name in prime_names.values()}
    early_warnings = []
    warning_threshold = 0.8  # fraction of target before warning

    for t, snapshot in enumerate(history):
        # Temporarily set space state to snapshot for aggregate computation
        for c, v in snapshot.items():
            space.set(c, v)
        for p, name in prime_names.items():
            agg = space.aggregate(p)
            aggregates[name].append(agg)
            tgt = target_errors.get(p, 1.0)
            if abs(agg) > warning_threshold * abs(tgt):
                early_warnings.append((t, name, agg))

    return {
        "history": history,
        "aggregates": aggregates,
        "early_warnings": early_warnings,
    }


# --- Example usage ---
if __name__ == "__main__":
    # Initialize errors for all 16 binary configurations
    error_init = {}
    rng = np.random.default_rng(1)
    for m in [0, 1]:
        for r in [0, 1]:
            for o in [0, 1]:
                for c in [0, 1]:
                    error_init[(m, r, o, c)] = rng.normal(0, 0.1)

    target_errors = {
        P_MACHINE:   0.0,   # drive machine aggregate to zero bias
        P_REAGENT:   0.0,
        P_OPERATOR:  0.0,
        P_PATIENT:   0.0,
    }

    results = run_lab_simulation(error_init, target_errors, T=300)
    print(f"Early warnings detected: {len(results['early_warnings'])}")
    for t, name, agg in results['early_warnings'][:5]:
        print(f"  t={t:3d}  mode={name:10s}  aggregate={agg:.4f}")
\end{lstlisting}

\subsection{Epidemiological SIR with Prime-Encoded Multiplicity}

\begin{lstlisting}[language=Python, caption={Prime-encoded SIR epidemiological model with recursive stability constraints.}]
"""
sir_multiplicity.py
Prime-encoded SIR epidemiological model.
Prime modes: p_h (host), p_p (pathogen strain), p_s (contact setting)
"""

import numpy as np
from dataclasses import dataclass
from typing import Dict, List, Tuple

P_HOST     = 2   # host type
P_PATHOGEN = 3   # pathogen strain
P_SETTING  = 5   # contact setting

@dataclass
class SIRConfig:
    """A prime-labeled epidemiological configuration."""
    host:     int  # 0 or 1
    pathogen: int  # 0 or 1
    setting:  int  # 0 or 1

    @property
    def label(self) -> int:
        n = 1
        if self.host:     n *= P_HOST
        if self.pathogen: n *= P_PATHOGEN
        if self.setting:  n *= P_SETTING
        return n


def make_state(n_configs: int) -> Dict[int, Dict[str, float]]:
    """State: {label: {'S': ..., 'I': ..., 'R': ...}}."""
    configs = []
    for h in [0, 1]:
        for p in [0, 1]:
            for s in [0, 1]:
                configs.append(SIRConfig(h, p, s))
    state = {}
    for cfg in configs:
        state[cfg.label] = {'S': 0.95, 'I': 0.05, 'R': 0.0}
    return state, configs


def sir_step(
    state: Dict[int, Dict[str, float]],
    configs: List[SIRConfig],
    beta_tensor: Dict[Tuple[int, int], float],
    gamma: float,
    constraints: Dict[int, Tuple[float, float]],  # prime -> (target_I, gain)
    dt: float = 0.5,
) -> Dict[int, Dict[str, float]]:
    """
    One discrete step of prime-encoded SIR dynamics with stability constraints.
    beta_tensor: {(label_i, label_j): beta_ij} -- transmission from j to i.
    constraints: {prime: (target_aggregate_I, gain)}.
    """
    new_state = {lbl: dict(vals) for lbl, vals in state.items()}

    # 1. Raw SIR dynamics
    for cfg in configs:
        lbl = cfg.label
        S, I, R = state[lbl]['S'], state[lbl]['I'], state[lbl]['R']
        # Force from all interacting configurations
        force = sum(
            beta_tensor.get((lbl, cfg2.label), 0.0) * state[cfg2.label]['I']
            for cfg2 in configs
        )
        dS = -force * S * dt
        dI = (force * S - gamma * I) * dt
        dR = gamma * I * dt
        new_state[lbl]['S'] = max(0.0, S + dS)
        new_state[lbl]['I'] = max(0.0, I + dI)
        new_state[lbl]['R'] = max(0.0, R + dR)
        # Normalize to preserve total population per config
        total = new_state[lbl]['S'] + new_state[lbl]['I'] + new_state[lbl]['R']
        if total > 0:
            for k in ['S', 'I', 'R']:
                new_state[lbl][k] /= total

    # 2. Prime-wise stability corrections on I
    for prime, (target_I, gain) in constraints.items():
        current_agg = sum(
            new_state[cfg.label]['I']
            for cfg in configs
            if cfg.label % prime == 0
        )
        correction = gain * (target_I - current_agg)
        relevant = [cfg for cfg in configs if cfg.label % prime == 0]
        if relevant:
            per_cfg = correction / len(relevant)
            for cfg in relevant:
                new_state[cfg.label]['I'] = max(
                    0.0, new_state[cfg.label]['I'] + per_cfg
                )

    return new_state


def run_sir_multiplicity(T: int = 200, dt: float = 0.5) -> dict:
    """Run prime-encoded SIR simulation and return results."""
    state, configs = make_state(8)

    # Heterogeneous transmission: higher for same pathogen strain
    beta_tensor = {}
    for c1 in configs:
        for c2 in configs:
            base = 0.3
            # Same pathogen strain boosts transmission
            if c1.pathogen == c2.pathogen == 1:
                base *= 1.8
            # Same setting boosts transmission
            if c1.setting == c2.setting == 1:
                base *= 1.4
            beta_tensor[(c1.label, c2.label)] = base

    gamma = 0.1
    # Target: drive total infected per pathogen-prime toward 0.05 (endemic)
    constraints = {
        P_HOST:     (0.10, 0.05),
        P_PATHOGEN: (0.08, 0.05),
        P_SETTING:  (0.12, 0.05),
    }

    history = []
    aggregate_I = {P_HOST: [], P_PATHOGEN: [], P_SETTING: []}
    early_warnings = []

    for t in range(T):
        history.append({lbl: dict(v) for lbl, v in state.items()})

        for prime, prime_name in [(P_HOST, 'host'),
                                  (P_PATHOGEN, 'pathogen'),
                                  (P_SETTING, 'setting')]:
            agg = sum(
                state[cfg.label]['I'] for cfg in configs
                if cfg.label % prime == 0
            )
            aggregate_I[prime].append(agg)
            target = constraints[prime][0]
            if agg > 1.5 * target:
                early_warnings.append((t, prime_name, agg))

        state = sir_step(state, configs, beta_tensor, gamma, constraints, dt)

    return {
        "history": history,
        "aggregate_I": aggregate_I,
        "early_warnings": early_warnings,
        "configs": configs,
    }


if __name__ == "__main__":
    results = run_sir_multiplicity(T=300)
    print("Early-warning events (first 5):")
    for t, name, val in results['early_warnings'][:5]:
        print(f"  t={t:3d}  prime={name:10s}  aggregate_I={val:.4f}")
    print("\nFinal aggregate infections by prime:")
    for prime, name in [(P_HOST, 'host'), (P_PATHOGEN, 'pathogen'),
                        (P_SETTING, 'setting')]:
        print(f"  {name}: {results['aggregate_I'][prime][-1]:.4f}")
\end{lstlisting}

\subsection{Population Dynamics: Prime-Encoded Lotka--Volterra}

\begin{lstlisting}[language=Python, caption={Prime-encoded Lotka--Volterra system with age- and patch-structured prime modes.}]
"""
lotka_volterra_multiplicity.py
Prime-encoded predator-prey dynamics.
Prime modes: p_s (species), p_a (age class), p_n (niche/patch)
"""

import numpy as np

P_SPECIES = 2   # species: 0=prey, 1=predator
P_AGE     = 3   # age class: 0=juvenile, 1=adult
P_NICHE   = 5   # niche/patch: 0=patch A, 1=patch B


def label(species, age, niche):
    n = 1
    if species: n *= P_SPECIES
    if age:     n *= P_AGE
    if niche:   n *= P_NICHE
    return n


def contains(lbl, prime):
    return lbl % prime == 0


def build_initial_state():
    """Build initial population tensor."""
    state = {}
    for s in [0, 1]:
        for a in [0, 1]:
            for n in [0, 1]:
                lbl = label(s, a, n)
                # Prey denser than predators, adults more than juveniles
                if s == 0:    # prey
                    state[lbl] = 10.0 if a == 1 else 4.0
                else:         # predator
                    state[lbl] = 2.0 if a == 1 else 0.8
    return state


def lv_step(
    state: dict,
    alpha: dict,       # {label: birth rate}
    beta:  dict,       # {(prey_lbl, pred_lbl): predation rate}
    death: dict,       # {label: death rate}
    constraints: dict, # {prime: (target_agg, gain)}
    dt: float = 0.05,
) -> dict:
    """One Euler step of prime-encoded LV with stability corrections."""
    new_state = {}
    all_labels = list(state.keys())

    for lbl in all_labels:
        P = state[lbl]
        birth_flux = alpha.get(lbl, 0.0) * P * dt
        death_flux = death.get(lbl, 0.0) * P * dt
        pred_flux  = sum(
            beta.get((lbl, pred_lbl), 0.0) * P * state[pred_lbl] * dt
            for pred_lbl in all_labels if contains(pred_lbl, P_SPECIES)
        )
        gain_flux  = sum(
            beta.get((prey_lbl, lbl), 0.0) * state[prey_lbl] * P * dt
            for prey_lbl in all_labels if not contains(prey_lbl, P_SPECIES)
        )
        new_state[lbl] = max(0.0, P + birth_flux - death_flux
                                    - pred_flux + gain_flux)

    # Prime-wise stability corrections
    for prime, (target, gain) in constraints.items():
        agg = sum(v for lbl, v in new_state.items() if contains(lbl, prime))
        correction = gain * (target - agg)
        relevant = [lbl for lbl in all_labels if contains(lbl, prime)]
        if relevant:
            per_lbl = correction / len(relevant)
            for lbl in relevant:
                new_state[lbl] = max(0.0, new_state[lbl] + per_lbl)

    return new_state


def run_lv_multiplicity(T: int = 500, dt: float = 0.05) -> dict:
    state = build_initial_state()
    all_labels = list(state.keys())

    # Birth rates higher for adults
    alpha = {lbl: (0.6 if not contains(lbl, P_SPECIES) else 0.0)
             for lbl in all_labels}
    # Predation: predator-adult on prey-adult highest
    beta = {}
    for prey_lbl in all_labels:
        for pred_lbl in all_labels:
            if not contains(prey_lbl, P_SPECIES) and contains(pred_lbl, P_SPECIES):
                rate = 0.05
                if contains(prey_lbl, P_AGE) and contains(pred_lbl, P_AGE):
                    rate *= 2.0  # adult-adult highest predation
                beta[(prey_lbl, pred_lbl)] = rate
    death = {lbl: (0.1 if contains(lbl, P_SPECIES) else 0.05)
             for lbl in all_labels}

    # Stability: drive prey aggregate toward 20, predator toward 5
    prey_labels  = [lbl for lbl in all_labels if not contains(lbl, P_SPECIES)]
    pred_labels  = [lbl for lbl in all_labels if contains(lbl, P_SPECIES)]
    constraints = {
        P_SPECIES: (5.0,  0.02),  # predator aggregate target
        P_AGE:     (15.0, 0.01),  # adult aggregate target
        P_NICHE:   (10.0, 0.01),  # patch aggregate target
    }

    history = []
    aggregates = {P_SPECIES: [], P_AGE: [], P_NICHE: []}
    early_warnings = []

    for t in range(T):
        history.append(dict(state))
        for prime in [P_SPECIES, P_AGE, P_NICHE]:
            agg = sum(v for lbl, v in state.items() if contains(lbl, prime))
            aggregates[prime].append(agg)
            target = constraints[prime][0]
            if agg > 3.0 * target:
                early_warnings.append((t, prime, agg))

        state = lv_step(state, alpha, beta, death, constraints, dt)

    return {
        "history": history,
        "aggregates": aggregates,
        "early_warnings": early_warnings,
    }


if __name__ == "__main__":
    results = run_lv_multiplicity(T=400)
    print("Final aggregates:")
    for prime, name in [(P_SPECIES,'predator'),(P_AGE,'adult'),(P_NICHE,'patch-A')]:
        final = results['aggregates'][prime][-1]
        print(f"  {name:12s}: {final:.3f}")
    print(f"Early warnings: {len(results['early_warnings'])}")
\end{lstlisting}

\newpage
% ═══════════════════════════════════════════════════════════════════════════════
\section{Validation Roadmap}
\label{sec:validation}
% ═══════════════════════════════════════════════════════════════════════════════

\subsection{Immediate (0--6 months)}

\begin{itemize}[leftmargin=*]
  \item \textbf{Clinical lab pilot}: Apply the prime-encoded QC model to a
    retrospective dataset from a high-complexity clinical laboratory. Define
    machine ($p_m$), reagent ($p_r$), and operator ($p_o$) primes. Compare
    early-warning performance against standard Shewhart/CUSUM charts.
  \item \textbf{Epidemic time-series}: Fit the prime-encoded SIR model to a
    published seasonal influenza dataset stratified by age group and strain.
    Evaluate peak timing and final-size prediction against standard
    age-structured SIR.
\end{itemize}

\subsection{Medium-term (6--18 months)}

\begin{itemize}[leftmargin=*]
  \item \textbf{Ecological dataset}: Apply the prime-encoded Lotka--Volterra
    model to long-term multi-species time series (e.g., the Hudson Bay
    lynx-hare dataset or the Plankton mesocosm experiments). Test for
    prime-channel structure in residual correlations.
  \item \textbf{Neural recording data}: Apply the prime-encoded neural model
    to Neuropixels or Allen Institute datasets. Define brain-region ($p_r$)
    and cell-type ($p_n$) primes. Test predictions on cross-region oscillation
    coherence and early-warning for seizure transitions.
  \item \textbf{Evolutionary time series}: Apply the prime-encoded evolutionary
    model to SARS-CoV-2 variant surveillance data (host type $\times$ strain
    $\times$ region primes). Test predictions on variant replacement dynamics.
\end{itemize}

\subsection{Long-term (18--36 months)}

\begin{itemize}[leftmargin=*]
  \item \textbf{Multi-omics systems biology}: Apply to time-series
    metabolomics/transcriptomics in a cell-perturbation dataset. Recover
    regulatory/metabolic interaction tensors and compare to Boolean/ODE
    network models.
  \item \textbf{Reaction--diffusion experiments}: Apply to spatially resolved
    concentration data from synthetic Turing-patterning systems or
    developmental biology datasets. Test constrained pattern spectrum
    predictions.
  \item \textbf{Cross-domain coupling}: Design a joint eco-epidemiological
    model sharing host-type primes between the population-dynamics and
    epidemiology sectors, and fit to a dataset with both ecological and
    epidemiological time series.
\end{itemize}

\newpage
% ═══════════════════════════════════════════════════════════════════════════════
\section{Predictions Summary Table}
\label{sec:predictions}
% ═══════════════════════════════════════════════════════════════════════════════

\begin{longtable}{p{3cm} p{4cm} p{4cm} p{3.5cm}}
\toprule
\textbf{Domain} & \textbf{Structured Correlations} &
\textbf{Emergent Patterns} & \textbf{Early Warning} \\
\midrule
\endfirsthead
\multicolumn{4}{c}{\textit{Continued from previous page}} \\
\toprule
\textbf{Domain} & \textbf{Structured Correlations} &
\textbf{Emergent Patterns} & \textbf{Early Warning} \\
\midrule
\endhead
\midrule
\multicolumn{4}{r}{\textit{Continued on next page}} \\
\endfoot
\bottomrule
\endlastfoot

Population Dynamics &
Abundance covaries across age/patch sharing a species prime &
Localized boom/crash in specific genotype-niche configs &
$\mathcal{P}_p \to$ carrying-capacity boundary flags extinction risk \\[6pt]

Epidemiology &
All host types sharing a strain prime rise/fall together &
Multi-phasic early growth due to strain $\times$ setting interactions &
$\mathcal{I}_p$ exceeding endemic level flags outbreak \\[6pt]

Neural Networks &
Co-oscillation across region-sharing configurations &
Structured connectivity motifs for prime-wise stability &
$\mathcal{N}_p$ near seizure threshold flags pathological transition \\[6pt]

Ecological Networks &
Energy-flow covariance along trophic-level primes &
Emergent cycles at prime-combination periods &
$\mathcal{E}_p$ approaching energy limit flags collapse \\[6pt]

Evolutionary Dynamics &
Allele frequency covariance across deme/environment primes &
Hidden epistatic motifs in genotype$\times$environment configs &
$\mathcal{G}_p$ approaching boundary flags sweep or diversity loss \\[6pt]

Reaction-Diffusion &
Co-patterning of species sharing compartment primes &
Constrained wavelengths vs.\ standard Turing &
$\mathcal{C}_p$ at threshold flags pattern-type transition \\[6pt]

Systems Biology &
Multi-omics covariance along gene-module/cell-type primes &
Emergent metabolic-regulatory circuits in composite configs &
$\mathcal{M}_p$ or $\mathcal{E}_p$ boundary flags disease transition \\[6pt]

Clinical Lab Systems &
Error correlation along machine $\times$ reagent prime pairs &
Structured non-i.i.d.\ error topology invisible to standard QC &
$\mathcal{E}_p$ boundary flags impending out-of-control event \\

\end{longtable}

\newpage
% ═══════════════════════════════════════════════════════════════════════════════
\section{Defensive Prior Art Disclosure}
\label{sec:defensive}
% ═══════════════════════════════════════════════════════════════════════════════

\begin{disclosurebox}
\textbf{Document ID}: MTF-BIOL-2026-001 \\
\textbf{Publication Date}: April 25, 2026 \\
\textbf{Disclosure Authority}: Multiplicity Foundation, Bridgeville, PA, USA \\[0.5em]
\textbf{Purpose}: To establish prior art as of the publication date and prevent
third parties from obtaining patent rights over the subject matter described
herein. This document is intentionally and irrevocably published into the
public domain.
\end{disclosurebox}

\subsection{Scope of Disclosed Subject Matter}

The following subject matter is disclosed as prior art:

\begin{enumerate}[leftmargin=*]
  \item \textbf{Prime-indexed tensor biological state spaces}: Any mathematical
    framework in which biological state variables are indexed by multisets of
    prime numbers, including all specific prime assignments to biological
    entities (species, age class, allele, brain region, metabolite, machine,
    reagent, etc.).

  \item \textbf{Recursive stability constraints on prime aggregates}: Any method
    in which the aggregate $\mathcal{X}_p(t) = \sum_{\bp \ni p} X_{\bp}(t)$
    is constrained by a function $\Phi_p$, applied to any biological domain
    including but not limited to epidemiology, ecology, neuroscience,
    evolutionary biology, systems biology, and clinical laboratory QC.

  \item \textbf{The multiplicity meta-equation}: The equation
    $H(t,\psi) \to M(t,\psi) T(t,G) + f(t,\psi) = \lambda(t)\psi$
    and any instantiation thereof in biological modeling contexts, including
    the specific operator assignments ($M$ = stability constraints;
    $T$ = structural networks; $f$ = exogenous drivers; $\lambda$ = global
    consistency).

  \item \textbf{Quantum-inspired superposition on multiplicity basis states}:
    Any representation $\Psi(t) = \sum_{\bp} a_{\bp}(t)\psi_{\bp}$ where
    $\{\psi_{\bp}\}$ are prime-factorized biological configurations.

  \item \textbf{Prime-channel early-warning signals}: Any method for
    detecting impending biological regime transitions (epidemic waves,
    population crashes, neural seizures, lab failures, evolutionary sweeps,
    pattern transitions, cellular phase shifts) by monitoring prime-aggregate
    quantities $\mathcal{X}_p(t)$ relative to stability-envelope boundaries.

  \item \textbf{All code implementations}: All Python modules disclosed in
    Section~\ref{sec:code}, including the data structures
    (\texttt{Configuration}, \texttt{MultiplicitySpace},
    \texttt{StabilityConstraint}, \texttt{MultiplicitySystem}) and the domain
    implementations (\texttt{lab\_qc\_multiplicity.py},
    \texttt{sir\_multiplicity.py}, \texttt{lotka\_volterra\_multiplicity.py}).

  \item \textbf{Domain-specific instantiations}: All eight domain models
    described in Section~\ref{sec:domains}, including specific prime-mode
    assignments, state tensor definitions, delta functions, and
    stability-function parameterizations.

  \item \textbf{Cross-domain coupling via shared primes}: Any method of
    coupling biological sectors (e.g., epidemiology and population dynamics)
    through shared prime modes and joint prime-aggregate constraints.
\end{enumerate}

\subsection{Prior Art Context}

This disclosure builds on and extends the following bodies of prior work,
none of which discloses the prime-indexed multiplicity framework described
herein:

\begin{itemize}[leftmargin=*]
  \item Standard tensor methods in biomedical data analysis (e.g.,
    Tucker/CP decomposition, tensor networks), which do not use prime
    factorization for index labeling or impose prime-wise stability
    constraints.
  \item Classical mathematical biology models (Lotka--Volterra, SIR,
    Hodgkin--Huxley, Turing, Fisher equation), which do not employ
    prime-labeled index sets or the meta-equation structure.
  \item Generic factored state-space frameworks (e.g., FACTS,
    state-space models) in machine learning and dynamical systems,
    which do not use prime factorization or biological multiplicity.
  \item Random-with-constraints modeling paradigms in high-dimensional
    biology, which impose statistical but not prime-algebraic constraints.
  \item Existing Multiplicity Theory formulations (pre-2026), which
    defined the general prime-indexed recursive framework but had not
    applied it to the eight specific mathematical biology domains with
    domain-specific prime modes, stability functions, and validation
    protocols as disclosed here.
\end{itemize}

\subsection{Legal Notice}

This document is published openly and is freely accessible to the public.
By virtue of this publication, all subject matter described herein
constitutes prior art under 35 U.S.C. \S 102 (United States) and
equivalent provisions in other jurisdictions. Any patent application filed
after April 25, 2026 on subject matter that is identical or substantially
similar to the content of this document is hereby pre-empted.

No rights are reserved by the Multiplicity Foundation on the specific
technical disclosures contained herein. The framework as disclosed is
available for unrestricted use, implementation, and further development
by any party.

\newpage
% ═══════════════════════════════════════════════════════════════════════════════
\section{Future Directions}
\label{sec:future}
% ═══════════════════════════════════════════════════════════════════════════════

\begin{enumerate}[leftmargin=*]
  \item \textbf{Formal algebraic theory}: Prove existence/uniqueness theorems
    for recursively stable states in general multiplicity spaces;
    characterize conditions under which the meta-equation eigenvalue problem
    has a discrete spectrum.

  \item \textbf{Computational scaling}: Develop efficient algorithms for
    fitting prime-constrained models to large biological datasets,
    including tensor-sparse representations and GPU-accelerated
    multiplicity-space operations.

  \item \textbf{Visualization tools}: Create interactive simulation
    environments for exploring prime-mode dynamics, stability envelopes,
    and cross-sector coupling in real time.

  \item \textbf{Higher-exponent models}: Extend beyond binary ($e \in \{0,1\}$)
    exponents to quantitative prime multiplicities, enabling continuous
    encoding of concentrations, abundances, and expression levels.

  \item \textbf{Non-commutative extension}: Investigate ordered or directed
    prime products to encode directional (e.g., temporal, causal) structure
    not captured by commutative multisets.

  \item \textbf{Integration with machine learning}: Embed the multiplicity
    framework as a structural prior in neural networks or Gaussian process
    models, enabling data-driven discovery of prime modes and interaction
    tensors from multi-scale biological datasets.

  \item \textbf{Environmental and synthetic biology applications}: Apply the
    framework to environmental sustainability modeling (species,
    nutrient, climate primes) and synthetic gene circuit design
    (promoter, regulator, chassis primes).
\end{enumerate}

% ═══════════════════════════════════════════════════════════════════════════════
\section{Conclusion}
\label{sec:conclusion}
% ═══════════════════════════════════════════════════════════════════════════════

\MT{} for Mathematical Biology provides a rigorous, unified language for
modeling the multi-scale, multi-factor structure of biological systems. By
replacing generic tensor indices with prime-factorized configuration labels,
imposing algebraically motivated recursive stability constraints, and grounding
quantum-inspired superpositions in an explicit multiplicity basis, the framework
transforms eight classical areas of mathematical biology from isolated,
domain-specific models into a coherent family of sector instantiations of a
single governing meta-equation.

The framework's three universal predictions---structured prime-channel
correlations, emergent cross-scale interaction motifs, and early-warning
signals at stability boundaries---are empirically testable across all domains
and provide immediate targets for experimental and computational validation.

This document serves a dual purpose: as a technical report establishing the
mathematical foundations of MT-MB, and as a defensive prior art publication
placing the full intellectual content in the public domain as of April 25, 2026.

\appendix
\section*{Mathematical Appendix: Operator Formulation and Norm Bounds}
\addcontentsline{toc}{section}{Mathematical Appendix: Operator Formulation and Norm Bounds}

\setcounter{section}{0}
\renewcommand{\thesection}{A.\arabic{section}}
\renewcommand{\thesubsection}{A.\arabic{subsection}}
\renewcommand{\thetheorem}{A.\arabic{theorem}}
\renewcommand{\thedefinition}{A.\arabic{definition}}
\renewcommand{\theproposition}{A.\arabic{proposition}}
\renewcommand{\thecorollary}{A.\arabic{corollary}}
\renewcommand{\theremark}{A.\arabic{remark}}

\section{Operator-Theoretic Reformulation of Multiplicity Dynamics}

In this appendix we make the implicit functional-analytic structure of the
Multiplicity Theory framework explicit. We identify the biological state space
with a Banach space, define the linear and nonlinear operators corresponding
to the multiplicity constraints and structural networks, and derive operator
norm bounds sufficient for existence, uniqueness, and stability of solutions
to the meta-equation
\[
  H(t,\psi(t)) \;\to\; M(t,\psi(t))\,T(t,G) + f(t,\psi(t)) = \lambda(t)\,\psi(t),
\]
thereby promoting the heuristic reasoning in the main text to explicit
mathematical results.

\subsection{State Spaces and Basic Operators}

Let $\Pcal$ be the set of primes representing biological modes and
$\Pcal^*$ the free commutative monoid of configurations. For a fixed
finite subset $\Pcal_0 \subset \Pcal$ and a chosen truncation of exponents
$e_i \in \{0,\dots,E_{\max}\}$ we obtain a finite configuration set
$\mathcal{C} \subset \Pcal^*$, which is sufficient for any practical
implementation and empirical application.

\begin{definition}[State space]
Let $N = |\mathcal{C}|$. We identify the multiplicity state with a vector
$X \in \mathbb{R}^N$ by choosing a fixed enumeration
$\mathcal{C} = \{\bp_1,\dots,\bp_N\}$ and setting
$X_i = X_{\bp_i}$ for $i=1,\dots,N$. We equip $\mathbb{R}^N$ with the norm
$\|X\| = \|X\|_2$ unless stated otherwise.
\end{definition}

\begin{definition}[Aggregate operator]
Let $\Pcal_0 = \{p_1,\dots,p_M\}$ be the set of active primes. Define the
linear operator $A: \mathbb{R}^N \to \mathbb{R}^M$ by
\[
  (AX)_j \;=\; \mathcal{X}_{p_j}(X)
  \;=\; \sum_{i: p_j \mid n(\bp_i)} X_i,
  \quad j=1,\dots,M.
\]
We call $A$ the \emph{aggregate operator}.
\end{definition}

By construction $A$ is a nonnegative matrix. Its operator norm in the Euclidean
norm satisfies the trivial bound
\[
  \|A\|_2 \;\leq\; \sqrt{MN},
\]
but in applications the sparsity and overlap structure of primes yield much
tighter bounds (e.g.\ each configuration contains at most $k$ primes, or each
prime divides at most $d$ configurations).

\begin{definition}[Multiplicity constraint operator]
Let $\Phi: \mathbb{R}^N \to \mathbb{R}^M$ be given componentwise by
$\Phi_j(X) = \Phi_{p_j}(X)$, where $\Phi_{p_j}$ is the prime-wise stability
function from the main text. Define the \emph{constraint operator}
$M_\Phi : \mathbb{R}^N \to \mathbb{R}^N$ by
\[
  M_\Phi(X)
  \;=\; A^\top W\,\Phi(X),
\]
where $W \in \mathbb{R}^{M \times M}$ is a diagonal weight matrix specifying
how constraint residuals are redistributed back to configuration space.
\end{definition}

\begin{remark}
The particular choice $W = \mathrm{diag}(w_1,\dots,w_M)$ with
$w_j = 1/|\{i : p_j \mid n(\bp_i)\}|$ corresponds exactly to the “distribute
equally over configurations containing prime $p_j$” rule implemented in the
reference code for clinical lab, SIR, and Lotka–Volterra modules.
\end{remark}

\begin{definition}[Structural transfer operator]
The structural interactions (predation, infection, synaptic transmission,
reaction fluxes, etc.) are encoded in a linear operator
$T : \mathbb{R}^N \to \mathbb{R}^N$, represented by an $N \times N$ matrix
$K = (K_{ij})$ with entries determined by the corresponding interaction
tensor (e.g.\ $K_{ij} = \beta_{\bp_i,\bp_j}$ for epidemiological coupling).
\end{definition}

\begin{definition}[External forcing]
Exogenous drives (seasonality, environmental forcing, lab reagent lot changes,
etc.) are represented by a time-dependent term $f(t) \in \mathbb{R}^N$, assumed
bounded and measurable in $t$.
\end{definition}

With these definitions, the discrete-time multiplicity dynamics in vector form
can be written as
\begin{equation}
  X(t+1) \;=\; X(t) + F\bigl(t,X(t)\bigr), \qquad
  F(t,X) \;=\; M_\Phi(X) + T X + f(t),
  \label{eq:disc-appendix}
\end{equation}
and the meta-equation becomes the requirement that $F(t,X)$ admit the
eigen-structure
\[
  F\bigl(t,\psi(t)\bigr) \;=\; \lambda(t)\,\psi(t)
\]
for the global multiplicity state $\psi(t)$.

\section{Lipschitz Bounds and Global Existence}

We first establish sufficient conditions for global existence and uniqueness of
solutions to the discrete-time system~\eqref{eq:disc-appendix} in the finite
dimensional setting.

\begin{assumption}[Lipschitz stability map]
There exists $L_\Phi \geq 0$ such that
\[
  \|\Phi(X) - \Phi(Y)\|_2 \;\leq\; L_\Phi\,\|X - Y\|_2
  \quad\text{for all }X,Y \in \mathbb{R}^N.
\]
\end{assumption}

This is satisfied in particular when each $\Phi_{p_j}$ is affine
$\Phi_{p_j}(X) = b_j - c_j^\top X$ with bounded $c_j$.

\begin{proposition}[Lipschitz bound for $M_\Phi$]
Under the above assumption,
\[
  \|M_\Phi(X) - M_\Phi(Y)\|_2
  \;\leq\; \|A^\top W\|_2\,L_\Phi\,\|X - Y\|_2.
\]
\end{proposition}

\begin{proof}
Using the operator norm inequality,
\[
  \|M_\Phi(X) - M_\Phi(Y)\|_2
  = \|A^\top W (\Phi(X) - \Phi(Y))\|_2
  \leq \|A^\top W\|_2 \,\|\Phi(X) - \Phi(Y)\|_2
  \leq \|A^\top W\|_2 L_\Phi \|X-Y\|_2.
\]
\end{proof}

\begin{assumption}[Bounded transfer operator]
The structural operator $T$ is bounded with
\[
  \|T\|_2 \;\leq\; L_T < \infty.
\]
\end{assumption}

This is automatic in finite dimension; $L_T$ may be taken as the spectral
norm of $K$.

\begin{proposition}[Lipschitz bound for full update]
Define $F(t,X) = M_\Phi(X) + T X + f(t)$ as in~\eqref{eq:disc-appendix}.
Assume $f(t)$ is bounded uniformly in $t$. Then $F$ is globally Lipschitz
in $X$ with
\[
  \|F(t,X) - F(t,Y)\|_2
  \;\leq\; L_F\,\|X-Y\|_2,\qquad
  L_F = \|A^\top W\|_2 L_\Phi + L_T.
\]
\end{proposition}

\begin{proof}
Immediate from the preceding bound and the fact that $f(t)$ cancels in the
difference:
\[
  \|F(t,X)-F(t,Y)\|_2
  \leq \|M_\Phi(X)-M_\Phi(Y)\|_2 + \|T(X-Y)\|_2
  \leq (\|A^\top W\|_2 L_\Phi + L_T)\,\|X-Y\|_2.
\]
\end{proof}

\begin{theorem}[Global existence and uniqueness]
For any initial condition $X(0) \in \mathbb{R}^N$ there exists a unique
solution $\{X(t)\}_{t\ge 0}$ to~\eqref{eq:disc-appendix}, defined for all
$t \in \mathbb{N}$.
\end{theorem}

\begin{proof}
The map $G_t(X) = X + F(t,X)$ is Lipschitz in $X$ for each $t$ with constant
$1+L_F$. Hence, by the Picard iteration principle for discrete-time systems,
given $X(t)$ there exists a unique $X(t+1)$ solving~\eqref{eq:disc-appendix}.
Inductively, the sequence is uniquely defined for all $t$.
\end{proof}

\section{Linearization and Prime-Channel Stability}

We now analyze the linearized dynamics around a recursively stable state and
establish bounds on the spectral radius of the linearization in terms of the
operator norms of $M_\Phi$ and $T$.

\begin{definition}[Recursively stable fixed point]
A state $X^\ast \in \mathbb{R}^N$ is \emph{recursively stable} if
\[
  F(t,X^\ast) \equiv 0 \quad\text{for all }t
\]
(hence $X(t) \equiv X^\ast$ is a solution), and the Jacobian matrix
$J = D_X F(t,X^\ast)$ has spectral radius $\rho(J) < 1$.
\end{definition}

\begin{remark}
The time-independence of $F(t,X^\ast)$ at a fixed point is guaranteed if the
forcing term satisfies $f(t) \equiv f^\ast$ such that
$M_\Phi(X^\ast) + T X^\ast + f^\ast = 0$ for all $t$ (e.g.\ stationary forcing,
or forcing that is absorbed into $X^\ast$ by construction).
\end{remark}

\begin{proposition}[Jacobian structure]
Suppose each $\Phi_{p_j}$ is differentiable. Let $J_\Phi = D_X \Phi(X^\ast)$
denote the $M \times N$ Jacobian of the stability map at $X^\ast$.
Then the Jacobian of $F$ at $X^\ast$ is
\[
  J = D_X F(t,X^\ast)
    = A^\top W\,J_\Phi + T.
\]
\end{proposition}

\begin{proof}
By linearity of $A^\top W$ and $T$,
\[
  D_X M_\Phi(X^\ast) = A^\top W\, D_X \Phi(X^\ast) = A^\top W J_\Phi,
\]
and $D_X(TX) = T$, while $D_X f(t) = 0$.
\end{proof}

\begin{theorem}[Sufficient spectral-norm stability condition]
If
\[
  \|A^\top W J_\Phi\|_2 + \|T\|_2 \;<\; 1,
\]
then $\rho(J) < 1$ and $X^\ast$ is exponentially stable: there exist
constants $C>0$ and $0<\kappa<1$ such that any solution with initial condition
$X(0)$ satisfies
\[
  \|X(t) - X^\ast\|_2 \;\leq\; C\,\kappa^t\,\|X(0)-X^\ast\|_2
  \quad\text{for all } t \in \mathbb{N}.
\]
\end{theorem}

\begin{proof}
By the triangle inequality for operator norms,
\[
  \|J\|_2
  = \|A^\top W J_\Phi + T\|_2
  \leq \|A^\top W J_\Phi\|_2 + \|T\|_2.
\]
If the right-hand side is $<1$ then $\|J\|_2<1$, hence all eigenvalues
satisfy $|\lambda| \le \|J\|_2 < 1$, so $\rho(J)<1$. Standard linear
discrete-time stability theory then yields exponential convergence to $X^\ast$.
\end{proof}

\begin{corollary}[Prime-channel stability]
If the stability map $J_\Phi$ is block-diagonal in the decomposition induced
by prime factors (i.e.\ configurations with disjoint prime supports do not
influence each other's stability) and the structural operator $T$ respects
this block structure, then the system decomposes into independent
\emph{prime-channel subsystems}. Each block $J^{(q)}$ indexed by a composite
prime product $q$ satisfies its own spectral-norm stability condition, and the
overall system is stable iff all prime-channel blocks are stable.
\end{corollary}

\section{Norm Bounds for Aggregates and Early-Warning Indicators}

We next obtain explicit bounds for the evolution of prime aggregates and show
how they lead to quantitative early-warning criteria.

\begin{proposition}[Aggregate sensitivity bound]
Let $X(t)$ satisfy~\eqref{eq:disc-appendix}. For a given prime $p_j$ and
time step $t\to t+1$,
\[
  |\mathcal{X}_{p_j}(t+1) - \mathcal{X}_{p_j}(t)|
  \;\leq\; \|A_{j,\cdot}\|_2 \,\|F(t,X(t))\|_2,
\]
where $A_{j,\cdot}$ denotes the $j$-th row of $A$.
\end{proposition}

\begin{proof}
We have
\[
  \mathcal{X}_{p_j}(t+1) - \mathcal{X}_{p_j}(t)
  = (AX(t+1) - AX(t))_j
  = (A F(t,X(t)))_j
  = A_{j,\cdot} F(t,X(t)),
\]
and the Cauchy–Schwarz inequality yields
\[
  |A_{j,\cdot} F(t,X(t))|
  \le \|A_{j,\cdot}\|_2 \,\|F(t,X(t))\|_2.\qedhere
\]
\end{proof}

\begin{corollary}[Uniform envelope bound]
Suppose $\|F(t,X(t))\|_2 \leq B$ uniformly in $t$ and $j$. Then
\[
  |\mathcal{X}_{p_j}(t) - \mathcal{X}_{p_j}(0)|
  \;\leq\; t\,\|A_{j,\cdot}\|_2\,B,
\]
so that the aggregate grows at most linearly with slope
$\|A_{j,\cdot}\|_2 B$.
\end{corollary}

In practice we want the aggregates to remain within a \emph{stability
envelope} $[L_j,U_j]$, with $L_j < U_j$ chosen according to biological
constraints (e.g.\ acceptable error baselines, endemic infection levels,
carrying capacities). The above bound can then be translated into a lower
bound on the time required for crossing the envelope, which in turn defines a
lead time for early-warning indicators.

\begin{proposition}[Lead time to envelope breach]
Assume $\mathcal{X}_{p_j}(0) \in (L_j,U_j)$, and $\mathcal{X}_{p_j}$ is
monotone in $t$ (or monotone over a given time window). If
$\|F(t,X(t))\|_2 \le B$ and $\|A_{j,\cdot}\|_2 \le C_j$, then the earliest
time $\tau$ such that $\mathcal{X}_{p_j}(\tau) \notin (L_j,U_j)$ satisfies
\[
  \tau \;\geq\; \frac{\min\{|\mathcal{X}_{p_j}(0)-L_j|,
                              |\mathcal{X}_{p_j}(0)-U_j|\}}{C_j B}.
\]
\end{proposition}

This lower bound on $\tau$ provides a quantitative guarantee of how long a
system can remain within safe bounds under worst-case change rates. In
empirical applications, one reverses the logic: estimates of $C_j$ and $B$
can be obtained from data, and observed changes in $\mathcal{X}_{p_j}$ can be
compared to these bounds to detect anomalously fast approaches to stability
boundaries.

\section{Continuous-Time Reaction--Diffusion Bounds}

For reaction--diffusion systems the state is a field
$C_{\bp}(x,t)$ on a spatial domain $\Omega \subset \mathbb{R}^d$, and the
equation
\begin{equation}
  \frac{\partial C_{\bp}}{\partial t}
  = D_{\bp}\,\nabla^2 C_{\bp} + F_{\bp}(C(\cdot,t)),
  \label{eq:RD-appendix}
\end{equation}
is naturally analyzed in the Hilbert space $L^2(\Omega)^N$.

\subsection{Semigroup Formulation}

Let $H = L^2(\Omega)^N$ with inner product
\[
  \langle C, U \rangle
  = \sum_{i=1}^N \int_\Omega C_{\bp_i}(x)\,U_{\bp_i}(x)\,dx.
\]
Define the linear diffusion operator
\[
  (L C)_{\bp_i}
  = D_{\bp_i}\,\nabla^2 C_{\bp_i},
\]
with appropriate boundary conditions (e.g.\ homogeneous Neumann or Dirichlet)
so that $L$ generates an analytic contraction semigroup
$e^{tL}$ on $H$. Under homogeneous Neumann conditions and $D_{\bp_i}\ge 0$,
$L$ is self-adjoint and nonpositive, hence $\|e^{tL}\|_{\mathcal{L}(H)} \le 1$.

\begin{assumption}[Lipschitz reaction term]
There exists $L_F \ge 0$ such that
\[
  \|F(C) - F(U)\|_H \;\leq\; L_F\,\|C-U\|_H
  \quad\text{for all }C,U \in H,
\]
where $F: H \to H$ is the nonlinear reaction functional collecting all
$F_{\bp}$.
\end{assumption}

\begin{theorem}[Mild solution and norm bound]
For any initial condition $C(0) \in H$ the reaction--diffusion system
\eqref{eq:RD-appendix} admits a unique mild solution given by
\[
  C(t) = e^{tL} C(0)
  + \int_0^t e^{(t-s)L} F\bigl(C(s)\bigr)\,ds.
\]
Moreover, if $F(0)=0$ and $L_F < \infty$, then
\[
  \|C(t)\|_H \;\leq\; \|C(0)\|_H\,e^{L_F t}.
\]
\end{theorem}

\begin{proof}
Standard semigroup theory for semilinear evolution equations applies. The
norm bound follows from
\[
  \|C(t)\|_H
  \le \|e^{tL} C(0)\|_H + \int_0^t \|e^{(t-s)L}\|\,\|F(C(s))\|_H\,ds
  \le \|C(0)\|_H + L_F \int_0^t \|C(s)\|_H\,ds,
\]
and Grönwall's inequality yields the stated estimate.
\end{proof}

\subsection{Prime-Aggregate Bounds in $L^2$}

Define the prime-aggregate functionals in continuous space by
\[
  \mathcal{C}_p(t)
  = \int_\Omega \sum_{\bp_i: p\mid n(\bp_i)} C_{\bp_i}(x,t)\,dx.
\]
Assuming $C_{\bp_i} \ge 0$ almost everywhere, we can bound their derivatives
by
\[
  \left|\frac{d}{dt}\,\mathcal{C}_p(t)\right|
  \le \sum_{\bp_i: p\mid n(\bp_i)}
    \int_\Omega \bigl| D_{\bp_i}\,\nabla^2 C_{\bp_i}
                     + F_{\bp_i}(C)\bigr|\,dx.
\]
Under homogeneous Neumann conditions the integral of the Laplacian vanishes,
so the bound simplifies to
\[
  \left|\frac{d}{dt}\,\mathcal{C}_p(t)\right|
  \le \sum_{\bp_i: p\mid n(\bp_i)} \|F_{\bp_i}(C)\|_{L^1(\Omega)}.
\]
If each $F_{\bp_i}$ is Lipschitz with constant $L_i$ and satisfies
$F_{\bp_i}(0)=0$, then
\[
  \|F_{\bp_i}(C)\|_{L^1}
  \le L_i\,\|C\|_{L^1} \le L_i\,|\Omega|^{1/2}\,\|C\|_H,
\]
so that
\[
  \left|\frac{d}{dt}\,\mathcal{C}_p(t)\right|
  \le \Bigl(\sum_{\bp_i: p\mid n(\bp_i)} L_i\Bigr)\,|\Omega|^{1/2}\,\|C(t)\|_H
  \le \Bigl(\sum_{\bp_i: p\mid n(\bp_i)} L_i\Bigr)\,|\Omega|^{1/2}\,\|C(0)\|_H e^{L_F t}.
\]
This yields explicit growth bounds for prime aggregates in reaction--diffusion
systems, which can be used in the same way as in the discrete-time case to
derive lead times for pattern transitions (e.g.\ from homogeneous states to
Turing patterns) under the multiplicity constraints.

\section{Summary of Key Norm Conditions}

For ease of reference we collect the main sufficient conditions derived in
this appendix:

\begin{enumerate}
  \item \textbf{Global well-posedness (discrete time)}:\\
    If the stability map $\Phi$ is Lipschitz with constant $L_\Phi$ and
    $\|T\|_2 \le L_T$, then the full update $F$ is Lipschitz with
    $L_F = \|A^\top W\|_2 L_\Phi + L_T$, and the system admits a unique global
    solution.

  \item \textbf{Linear stability near fixed points}:\\
    If $X^\ast$ satisfies $F(t,X^\ast)=0$ and
    $\|A^\top W J_\Phi\|_2 + \|T\|_2 < 1$, where $J_\Phi = D_X \Phi(X^\ast)$,
    then $X^\ast$ is exponentially stable with rate bounded by $\|J\|_2$.

  \item \textbf{Prime-channel decoupling}:\\
    If $J_\Phi$ and $T$ are block-diagonal in the prime-channel decomposition,
    the system splits into independent sub-systems whose stability can be
    analyzed separately.

  \item \textbf{Aggregate growth bounds}:\\
    Aggregate changes satisfy
    $|\mathcal{X}_{p_j}(t+1) - \mathcal{X}_{p_j}(t)|
     \le \|A_{j,\cdot}\|_2\|F(t,X(t))\|_2$, leading to linear-in-time envelope
    bounds and explicit lead-time estimates for crossing stability thresholds.

  \item \textbf{Reaction--diffusion norm bounds}:\\
    With Lipschitz reaction term $F$ and diffusion generator $L$, the
    semilinear PDE admits a unique mild solution with
    $\|C(t)\|_H \le \|C(0)\|_H e^{L_F t}$, and aggregate functionals
    $\mathcal{C}_p(t)$ satisfy explicit derivative bounds in terms of the
    Lipschitz constants of $F_{\bp}$ and the $L^2$ norm of $C$.
\end{enumerate}

Together, these results supply a rigorous backbone for the heuristic notions
of ``recursive stability'', ``prime channels'', and ``early-warning signals''
used throughout the main text, and provide concrete operator-norm conditions
under which the multiplicity framework yields globally well-behaved and
predictive models across all domains considered.

\nocite{*}
\bibliographystyle{unsrt}  
\bibliography{references}
\end{document}